% Options for packages loaded elsewhere
\PassOptionsToPackage{unicode}{hyperref}
\PassOptionsToPackage{hyphens}{url}
%
\documentclass[
  11pt,
]{article}
\usepackage{amsmath,amssymb}
\usepackage{iftex}
\ifPDFTeX
  \usepackage[T1]{fontenc}
  \usepackage[utf8]{inputenc}
  \usepackage{textcomp} % provide euro and other symbols
\else % if luatex or xetex
  \usepackage{unicode-math} % this also loads fontspec
  \defaultfontfeatures{Scale=MatchLowercase}
  \defaultfontfeatures[\rmfamily]{Ligatures=TeX,Scale=1}
\fi
\usepackage{lmodern}
\ifPDFTeX\else
  % xetex/luatex font selection
\fi
% Use upquote if available, for straight quotes in verbatim environments
\IfFileExists{upquote.sty}{\usepackage{upquote}}{}
\IfFileExists{microtype.sty}{% use microtype if available
  \usepackage[]{microtype}
  \UseMicrotypeSet[protrusion]{basicmath} % disable protrusion for tt fonts
}{}
\makeatletter
\@ifundefined{KOMAClassName}{% if non-KOMA class
  \IfFileExists{parskip.sty}{%
    \usepackage{parskip}
  }{% else
    \setlength{\parindent}{0pt}
    \setlength{\parskip}{6pt plus 2pt minus 1pt}}
}{% if KOMA class
  \KOMAoptions{parskip=half}}
\makeatother
\usepackage{xcolor}
\usepackage[margin=1in]{geometry}
\usepackage{longtable,booktabs,array}
\usepackage{calc} % for calculating minipage widths
% Correct order of tables after \paragraph or \subparagraph
\usepackage{etoolbox}
\makeatletter
\patchcmd\longtable{\par}{\if@noskipsec\mbox{}\fi\par}{}{}
\makeatother
% Allow footnotes in longtable head/foot
\IfFileExists{footnotehyper.sty}{\usepackage{footnotehyper}}{\usepackage{footnote}}
\makesavenoteenv{longtable}
\setlength{\emergencystretch}{3em} % prevent overfull lines
\providecommand{\tightlist}{%
  \setlength{\itemsep}{0pt}\setlength{\parskip}{0pt}}
\setcounter{secnumdepth}{-\maxdimen} % remove section numbering
\usepackage{amsmath,amssymb,amsthm,booktabs,longtable,microtype,xurl}
\emergencystretch=6em\sloppy
\ifLuaTeX
  \usepackage{selnolig}  % disable illegal ligatures
\fi
\usepackage{bookmark}
\IfFileExists{xurl.sty}{\usepackage{xurl}}{} % add URL line breaks if available
\urlstyle{same}
\hypersetup{
  pdftitle={A forbidden-subgraph theorem for graph-state monogamy},
  pdfauthor={Yongxian Zhang},
  hidelinks,
  pdfcreator={LaTeX via pandoc}}

\title{A forbidden-subgraph theorem for graph-state monogamy}
\author{Yongxian Zhang\thanks{School of Computer Science and Engineering,
South China University of Technology, Guangzhou, China.
Email: \href{mailto:mxymmxym1@gmail.com}{mxymmxym1@gmail.com}.
ORCID: \href{https://orcid.org/0009-0000-3864-3536}{0009-0000-3864-3536}.}}
\date{9 October 2026}

\begin{document}
\maketitle

\subsection{Abstract}\label{abstract}

We prove that positive tripartite information for a finite qubit graph
state forces a claw in some locally equivalent graph, establishing
Fuentes--Keeler--Munizzi--Pollack Conjecture 1. We classify all
claw-vertex-minor-free graphs: each connected component is locally
equivalent to one of six fixed graphs. A complete certificate for all
120 one-vertex extensions closes an induction over arbitrary graph
order. The graph-state entropy identity and exact cut profiles then give
the monogamy implication. Every connected graph on at least seven
vertices has a claw vertex-minor, and six is insufficient. Two
independently implemented exact checkers verify the finite proof data;
no complete Lean formalization is claimed.

\subsection{1. Exact statements}\label{exact-statements}

All graphs below are finite, simple, and undirected. A \textbf{claw} is
the four-vertex star \texttt{K\_\{1,3\}}. Local complementation at
\texttt{v}, denoted \texttt{G*v}, toggles edges between distinct
neighbors of \texttt{v} and changes no other edges. A vertex-minor is an
induced subgraph of a graph obtained by a finite sequence of local
complementations. This agrees with allowing local complementations and
vertex deletions in any order.

Write \texttt{R\_1=K\_1}, \texttt{R\_2=K\_2}, \texttt{R\_3=P\_3},
\texttt{R\_4=P\_4}, \texttt{R\_5=C\_5}, and \texttt{R\_6=W\_5}, where
\textbf{\texttt{W\_5} here has six vertices}, a five-cycle and its
universal hub. Equivalence to a representative permits graph isomorphism
as well as local complementation.

\textbf{Theorem 1 (complete component classification).} A graph has no
claw vertex-minor if and only if each nonempty connected component is
locally equivalent, up to isomorphism, to one of \texttt{R\_1,...,R\_6}.

\textbf{Theorem 2 (forbidden-subgraph implication).} For every qubit
graph state and every three pairwise disjoint subsystems \texttt{A,B,C},
if

{[}I\_3(A:B:C)=S(A)+S(B)+S(C)-S(AB)-S(AC)-S(BC)+S(ABC)\textgreater0.{]}

then the graph has a locally equivalent representative containing an
\textbf{induced} claw. Here \texttt{S} is von Neumann entropy in bits.
This proves the implication in Fuentes--Keeler--Munizzi--Pollack,
arXiv:2511.19585v1, Conjecture 1, without a bound on graph order.

\textbf{Corollary 3 (sharp connected-size threshold).} Every connected
graph on at least seven vertices has a claw vertex-minor. Six does not
suffice, because \texttt{W\_5} has no claw vertex-minor. For the
positive direction one can use at most eight local complementations,
whose pivot vertices lie in a connected set of at most seven vertices,
and then delete unwanted vertices. In the graph-state interpretation,
local Clifford gates followed by local computational-basis measurements
and Pauli corrections extract a four-qubit GHZ state on some four
vertices. This is not a claim about four prescribed vertices or about
general non-Clifford local operations.

The proof has a small finite certificate, \textbf{not} a finite-size
extrapolation. The sole structural finite step checks all 120 nonempty
one-vertex extensions of the six representatives. A connected-graph
induction proves that these cases cover arbitrary order.

\subsection{2. Local-complementation
lemmas}\label{local-complementation-lemmas}

\textbf{Restriction.} If \texttt{v} belongs to a vertex subset
\texttt{U}, then \texttt{(G*v){[}U{]}\ =\ (G{[}U{]})*v}. For two
vertices in \texttt{U}, whether their edge is toggled depends only on
their adjacencies to \texttt{v}, all of which are in \texttt{G{[}U{]}}.
Repeating this identity lifts any local-complementation sequence on an
induced subgraph to the full graph, with precisely the same restricted
result.

Consequently, vertex deletions can be postponed until the end of a
sequence: continue applying all later local complementations on the
original vertex set, and finally restrict to the vertices that were not
deleted. At each step the restriction agrees. In particular, having a
claw vertex-minor is inherited upward from induced subgraphs and is
invariant under local equivalence and isomorphism.

\textbf{Connectedness.} Local complementation preserves each connected
component, including its vertex set. It creates no edge between distinct
components. Any deleted edge \texttt{xy} had both ends adjacent to the
pivot \texttt{v}, so it is replaced, as a connectivity witness, by the
path \texttt{x-v-y}; these incident edges are unchanged. The inverse
operation is the same local complementation.

\textbf{Removable vertex.} A finite connected graph with at least two
vertices has a vertex whose deletion leaves a connected graph. Take a
spanning tree and delete any leaf of that tree; the remaining tree spans
all remaining vertices.

\textbf{Relabeling.} Relabeling a graph transports a local
complementation at \texttt{v} to local complementation at the relabeled
vertex. Therefore no relabeling is an extra physical or graph operation
in a lifted witness: it only changes how the next witness's vertex names
are interpreted.

\subsection{3. The finite extension
certificate}\label{the-finite-extension-certificate}

Represent a labeled graph on \texttt{0,...,m-1} by the integer adjacency
rows \texttt{a\_i}, whose \texttt{j}-th bit is one exactly when
\texttt{i} is adjacent to \texttt{j}. The six fixed representatives are:

\begin{longtable}[]{@{}lll@{}}
\toprule\noalign{}
m & Representative & Adjacency rows \\
\midrule\noalign{}
\endhead
\bottomrule\noalign{}
\endlastfoot
1 & K\_1 & 0 \\
2 & K\_2 & 2, 1 \\
3 & P\_3 & 2, 5, 2 \\
4 & P\_4 & 2, 5, 10, 4 \\
5 & C\_5 & 18, 5, 10, 20, 9 \\
6 & W\_5 & 50, 37, 42, 52, 41, 31 \\
\end{longtable}

For every \texttt{m=1,...,6} and
\texttt{1\ \textless{}=\ b\ \textless{}\ 2\^{}m}, add vertex \texttt{m}
with adjacency set given by the binary digits of \texttt{b}. These are
\textbf{all} nonempty attachment sets. The corresponding entry of
\texttt{certificate.json} supplies either:

\begin{itemize}
\tightlist
\item
  a list of local-complementation vertices and an ordered tuple
  \texttt{(center,leaf\_1,leaf\_2,leaf\_3)} inducing a claw afterward;
  or
\item
  a list of local-complementation vertices and an explicit bijection
  from \texttt{R\_\{m+1\}} to the entire resulting graph.
\end{itemize}

There is no latter outcome at \texttt{m=6}. All lists have length at
most three.

\begin{longtable}[]{@{}
  >{\raggedright\arraybackslash}p{(\columnwidth - 10\tabcolsep) * \real{0.1304}}
  >{\raggedleft\arraybackslash}p{(\columnwidth - 10\tabcolsep) * \real{0.1739}}
  >{\raggedleft\arraybackslash}p{(\columnwidth - 10\tabcolsep) * \real{0.1739}}
  >{\raggedleft\arraybackslash}p{(\columnwidth - 10\tabcolsep) * \real{0.1739}}
  >{\raggedleft\arraybackslash}p{(\columnwidth - 10\tabcolsep) * \real{0.1739}}
  >{\raggedleft\arraybackslash}p{(\columnwidth - 10\tabcolsep) * \real{0.1739}}@{}}
\toprule\noalign{}
\begin{minipage}[b]{\linewidth}\raggedright
Base order
\end{minipage} & \begin{minipage}[b]{\linewidth}\raggedleft
All attachments
\end{minipage} & \begin{minipage}[b]{\linewidth}\raggedleft
Claw endpoints
\end{minipage} & \begin{minipage}[b]{\linewidth}\raggedleft
R\_(m+1) endpoints
\end{minipage} & \begin{minipage}[b]{\linewidth}\raggedleft
Maximum LC length to a claw
\end{minipage} & \begin{minipage}[b]{\linewidth}\raggedleft
Maximum LC length to R\_(m+1)
\end{minipage} \\
\midrule\noalign{}
\endhead
\bottomrule\noalign{}
\endlastfoot
1 & 1 & 0 & 1 & --- & 0 \\
2 & 3 & 0 & 3 & --- & 1 \\
3 & 7 & 1 & 6 & 0 & 3 \\
4 & 15 & 11 & 4 & 2 & 2 \\
5 & 31 & 30 & 1 & 1 & 0 \\
6 & 63 & 63 & 0 & 2 & --- \\
Total & 120 & 105 & 15 & 2 & 3 \\
\end{longtable}

This table's role is fully specified by the following verifier,
expressed mathematically. For each pair \texttt{(m,b)} in the indicated
range, start with the fixed extension. For each pivot \texttt{v}, and
each row \texttt{i}, replace \texttt{a\_i} by
\texttt{a\_i\ XOR\ (a\_v\ XOR\ 2\^{}i)} if the \texttt{i}-th bit of
\texttt{a\_v} is one; otherwise retain \texttt{a\_i}. Use the old pivot
row throughout a step. At a claw endpoint verify all three center--leaf
adjacencies and all three leaf--leaf nonadjacencies. At a representative
endpoint verify that the supplied list is a permutation and compare
every matrix entry under that permutation. Finally require each pair
\texttt{(m,b)} exactly once. The scripts use explicit exceptions rather
than assertions; optimized Python does not remove checks.

\texttt{check\_certificate.py} implements this rule with binary rows.
\texttt{verify\_independent.py} implements local complementation using
symmetric differences of edge sets, generates attachments as subsets
rather than integer intervals, and compares endpoint edge sets. Neither
imports the generator or the other verifier. The witness generator is
optional and is not a trusted proof input.

\subsection{4. Proof of the classification: the unbounded
step}\label{proof-of-the-classification-the-unbounded-step}

Suppose a connected nonempty graph \texttt{G} has no claw vertex-minor.
Induct on its number \texttt{n} of vertices. The case \texttt{n=1} is
\texttt{R\_1}. For \texttt{n\textgreater{}=2}, choose a removable vertex
\texttt{w}, and put \texttt{H=G-w}. The graph \texttt{H} is connected
and has no claw vertex-minor. By induction it is locally equivalent, up
to isomorphism, to \texttt{R\_m} for some
\texttt{1\textless{}=m\textless{}=6}; necessarily \texttt{m=n-1}.

Lift the sequence for \texttt{H} to \texttt{G} and relabel the remaining
vertices as \texttt{0,...,m-1}, with \texttt{w} labeled \texttt{m}. The
resulting full graph is connected, so the attachment set of \texttt{w}
to \texttt{R\_m} is nonempty. It is exactly one of the 120 certified
extensions. A claw endpoint would contradict the hypothesis. Hence the
endpoint is \texttt{R\_\{m+1\}}, proving the induction. In particular
\texttt{m=6} is impossible, and no connected graph of order at least
seven avoids a claw vertex-minor.

The converse requires showing that none of the six representatives has a
claw in its LC orbit. We give a short spectral proof in Section 6. An
independent finite negative certificate is also supplied:
\texttt{certificate.json} contains LC-closed claw-free sets of sizes
\texttt{1,1,4,11,132,132}, each containing its representative. Closure
entails that every graph in the entire LC orbit is in the certified set.
The independent verifier regenerates the full orbits without using their
supplied contents and obtains exactly the same sets. There are 281
states and 1511 one-step closure checks in total.

Local complementation acts separately on components. An induced claw is
connected, so it cannot use different components. The connected
classification therefore gives Theorem 1 for arbitrary graphs, including
the empty graph, for which the assertion is vacuous.

For Corollary 3, take a connected induced subgraph on seven vertices,
obtained by growing a connected set one vertex at a time. Apply the
extension normalization to this graph. The maximum costs for successful
normalizations through orders two to six are \texttt{0,1,3,2,0}, and the
final claw witness costs at most two. The total is at most eight. An
earlier claw outcome costs no more. Lift the sequence to the original
graph before deleting the unwanted vertices. The negative six-vertex
example follows from the converse just proved.

\subsection{5. The graph-state entropy identity, with
proof}\label{the-graph-state-entropy-identity-with-proof}

For a graph \texttt{G=(V,E)} on \texttt{n} vertices define

\texttt{\textbar{}G\textgreater{}\ =\ 2\^{}(-n/2)\ sum\_(x\ in\ F\_2\^{}V)\ (-1)\^{}(sum\_\{uv\ in\ E\}\ x\_u\ x\_v)\ \textbar{}x\textgreater{}}.

For \texttt{X\ subset\ V}, let
\texttt{B\ =\ Gamma{[}X,V\textbackslash{}X{]}}, and write
\texttt{r\_G(X)=rank\_(F\_2)(B)}. Then

\texttt{S(X)\ =\ r\_G(X)}.

To see this, write the exponent as
\texttt{q\_X(x)+q\_Y(y)+x\^{}T\ B\ y}. The partial trace has entries

\texttt{rho\_X(x,x\textquotesingle{})\ =\ 2\^{}(-\textbar{}X\textbar{})\ (-1)\^{}(q\_X(x)+q\_X(x\textquotesingle{}))\ 1\_\{B\^{}T(x+x\textquotesingle{})=0\}}.

Indeed, the sum of a binary additive character over \texttt{F\_2\^{}Y}
is zero unless that character is trivial. Conjugating by the diagonal
matrix with entries \texttt{(-1)\^{}\{q\_X(x)\}} removes the signs. If
\texttt{K=ker(B\^{}T)}, the remaining matrix has one all-ones block for
each coset of \texttt{K}, each multiplied by
\texttt{2\^{}(-\textbar{}X\textbar{})}. Every block has size
\texttt{2\^{}(\textbar{}X\textbar{}-r)}. Thus the nonzero eigenvalues
are \texttt{2\^{}(-r)} with multiplicity \texttt{2\^{}r}, proving the
identity, also for empty subsystems and zero rank.

Cut rank is invariant under local complementation. If the pivot
\texttt{v} lies in \texttt{X}, its row in \texttt{B} stays fixed, and
every row indexed by a neighbor of \texttt{v} in \texttt{X} acquires
that fixed row. These elementary binary row additions preserve rank. If
the pivot is outside \texttt{X}, apply the analogous column operations.
Symmetry of the adjacency matrix gives
\texttt{r\_G(X)=r\_G(V\textbackslash{}X)}.

For completeness, the graph operation really is a local Clifford
operation. Let \texttt{K\_u=X\_u\ product\_\{w\ in\ N(u)\}\ Z\_w} be the
graph stabilizers and put

\texttt{U\_v\ =\ exp(-i\ pi\ X\_v/4)\ product\_\{u\ in\ N(v)\}\ exp(i\ pi\ Z\_u/4)}.

Conjugation preserves \texttt{K\_v}; it maps \texttt{K\_u} for a
neighbor \texttt{u} to
\texttt{K\textquotesingle{}\_u\ K\textquotesingle{}\_v}, where primes
denote the stabilizers of \texttt{G*v}, and fixes the generators of
nonneighbors. These identities follow from \texttt{X\ Z=-iY},
\texttt{Z\ X=iY} and quarter-turn Pauli conjugation. Both sets generate
the same stabilizer group after the change of generators. The
independent graph stabilizers have a unique joint +1 state, so
\texttt{U\_v\textbar{}G\textgreater{}} and
\texttt{\textbar{}G*v\textgreater{}} differ only by a global phase. This
is the standard local-complementation rule; see Van den
Nest--Dehaene--De Moor (2004).

\subsection{6. Exact cut profiles and exclusion of the
claw}\label{exact-cut-profiles-and-exclusion-of-the-claw}

For \texttt{P\_4}, every one-vertex cut has rank one. The three
unordered two-versus-two cuts have ranks \texttt{1,2,2}: representatives
are \texttt{\{0,1\}}, \texttt{\{0,2\}}, \texttt{\{0,3\}}. The
corresponding matrices, with complement columns in increasing order, are
respectively

\texttt{{[}{[}0,0{]},{[}1,0{]}{]},\ {[}{[}1,0{]},{[}1,1{]}{]},\ {[}{[}1,0{]},{[}0,1{]}{]}}.

For \texttt{C\_5}, every pair has cut rank two. By the cycle symmetries
it suffices to check an adjacent pair and a nonadjacent pair, for
example \texttt{\{0,1\}} and \texttt{\{0,2\}}. Their matrices are

\texttt{{[}{[}0,0,1{]},{[}1,0,0{]}{]}} and
\texttt{{[}{[}1,0,1{]},{[}1,1,0{]}{]}}.

Together with nonzero single-vertex cuts and complementation, this
proves

\texttt{r\_(C\_5)(X)=min(\textbar{}X\textbar{},5-\textbar{}X\textbar{})}.

For \texttt{W\_5}, every three-vertex cut has rank three. Every such cut
has one side consisting of the hub and two cycle vertices. Up to the
cycle symmetries, choose \texttt{\{5,0,1\}} or \texttt{\{5,0,2\}}; with
hub row first, the matrices are

\texttt{{[}{[}1,1,1{]},{[}0,0,1{]},{[}1,0,0{]}{]}} and
\texttt{{[}{[}1,1,1{]},{[}1,0,1{]},{[}1,1,0{]}{]}}.

Both have odd determinant. Full row rank of these three-vertex cuts
implies full row rank for smaller subsets: extend the subset to three
vertices and restrict the independent rows of that three-vertex cut.
Complementation gives

\texttt{r\_(W\_5)(X)=min(\textbar{}X\textbar{},6-\textbar{}X\textbar{})}.

The first three representatives cannot contain a four-vertex induced
subgraph. A graph locally equivalent to \texttt{P\_4} cannot be a claw
because a claw has rank one at every two-versus-two cut, whereas
\texttt{P\_4} does not.

Now consider any graph locally equivalent to \texttt{C\_5}. All its
two-vertex cut ranks are two. It consequently has minimum degree at
least two: if a vertex had degree zero or one, place it with its
possible neighbor in a two-vertex subset; its cut row is zero. If an
induced claw existed, each of its three leaves would have to be adjacent
to the sole vertex outside the claw. Any two leaves would then have
identical cut rows, contradicting rank two.

Similarly, every graph locally equivalent to \texttt{W\_5} has minimum
degree at least three: a vertex of degree at most two, together with its
neighbors and any needed padding, forms a three-vertex subset with a
zero cut row. In an induced claw, every leaf would therefore have to be
adjacent to both outside vertices. Again two leaves have identical cut
rows, contradicting the rank-two profile. This proves the converse in
Theorem 1 without relying on an orbit enumeration.

\subsection{7. All partitions and the monogamy
theorem}\label{all-partitions-and-the-monogamy-theorem}

Let \texttt{D=V\textbackslash{}(A\ union\ B\ union\ C)}, and use purity
to write

\texttt{I\_3\ =\ r(A)+r(B)+r(C)+r(D)-r(AB)-r(AC)-r(AD)}.

This expression is symmetric in the four blocks, because complementary
cuts have equal rank. If any block is empty, the expression is zero.
Thus components on at most three vertices always contribute zero.

For \texttt{P\_4}, a four-nonempty-block partition has four singleton
blocks. The value is \texttt{4-(1+2+2)=-1}. For \texttt{C\_5}, the only
possible positive block-size profile is \texttt{(2,1,1,1)}, giving
\texttt{5-(2+2+2)=-1}. For \texttt{W\_5}, the positive profiles are
\texttt{(3,1,1,1)} and \texttt{(2,2,1,1)}, giving respectively
\texttt{6-(2+2+2)=0} and \texttt{6-(2+3+3)=-2}.

For a disjoint union, every cut matrix is block diagonal after row and
column permutations, so ranks, and hence \texttt{I\_3}, add over
components. Therefore a graph with no claw vertex-minor satisfies every
MMI inequality. Taking the contrapositive proves Theorem 2.

More precisely, for any claw-vertex-minor-free graph and any four-block
partition,

\texttt{I\_3\ =\ -N\_4\ -\ N\_5\ -\ 2\ N\_6},

where \texttt{N\_4} counts four-vertex components meeting all four
blocks, \texttt{N\_5} counts five-vertex components meeting all four
blocks, and \texttt{N\_6} counts six-vertex components whose block-size
profile is \texttt{(2,2,1,1)}. This also completely describes all
equality cases \textbf{within the classified claw-vertex-minor-free
class}. It is not an equality classification of all graph-state MMI
instances.

\subsection{8. Scope, provenance, and independent
verification}\label{scope-provenance-and-independent-verification}

The original conjecture is from Fuentes, Keeler, Munizzi and Pollack,
\emph{Monogamy of Mutual Information in Graph States},
arXiv:2511.19585v1, Conjecture 1 (printed p.30). The source discusses a
necessary induced four-star somewhere in the local-Clifford orbit, not
an induced star in every representative. We prove exactly this
implication. Its additional generalized-star partition statement follows
explicitly: take I, J and K to be the three singleton leaves of the
induced claw and take C to be every other vertex, including the center.
The three leaf blocks have no edges between them and each is attached to
C, as required by the source definition.

The converse is not claimed. Claw vertex-minors concern what can be
obtained after deletions/measurements; those operations do not preserve
the original entropy vector. Nor do these results characterize all
MMI-satisfying graph states, handle arbitrary mixed or non-stabilizer
states, extend to odd-prime qudit graphs, or solve other holographic
entropy inequalities.

The finite LC and cut-rank tables for small graph states overlap the
established graph-state classification literature; those small states
themselves are not new. Our proof uses a complete extension induction
and explicit small witnesses rather than assuming an unrestricted
conclusion from a finite catalog. A bounded literature search did not
locate this exact implication's resolution, but this is \textbf{not}
evidence sufficient to assert historical priority. The separate 2026
paper on edgeless vertex-minor Ramsey numbers concerns a different
target graph and is not an input to this proof.

Both exact checkers pass all 120 extensions. They also check all 5460
assignments of representative vertices to four labeled blocks, including
empty blocks. The independent checker recomputes every orbit and
independently constructs the graph-state amplitudes: for all 126
representative cuts it verifies, by integer matrix multiplication, the
reduced-density numerator identity \texttt{M\^{}2=2\^{}(n-r)\ M} and
\texttt{tr(M)=2\^{}n}. These finite quantum checks supplement, rather
than replace, the all-order entropy proof in Section 5.

This is AI-assisted research. No external funding was received. No
external human peer review is claimed. Lean coverage, if supplied in a
later companion, must be stated separately; the present theorem is a
complete written proof with exact independently replayable finite
certificates, not a claim of complete Lean verification.

\small\subsection{References}\label{references}

\begin{enumerate}
\def\labelenumi{\arabic{enumi}.}
\tightlist
\item
  J. Fuentes, C. Keeler, W. Munizzi and J. Pollack, \emph{Monogamy of
  Mutual Information in Graph States}, arXiv:2511.19585v1 (2025),
  especially Conjecture 1.
\item
  M. Van den Nest, J. Dehaene and B. De Moor, \emph{Graphical
  description of the action of local Clifford transformations on graph
  states}, Physical Review A \textbf{69}, 022316 (2004), DOI
  10.1103/PhysRevA.69.022316.
\item
  M. Hein, J. Eisert and H. J. Briegel, \emph{Multiparty entanglement in
  graph states}, Physical Review A \textbf{69}, 062311 (2004), DOI
  10.1103/PhysRevA.69.062311.
\item
  J. de Jong, F. Hahn, N. Tcholtchev, M. Hauswirth and A. Pappa,
  \emph{Extracting GHZ states from linear cluster states}, Physical
  Review Research \textbf{6}, 013330 (2024), DOI
  10.1103/PhysRevResearch.6.013330. This is prior work on the path
  family, not an input to our arbitrary-connected-graph classification.
\item
  J. H. Bae, \emph{Vertex-minor Ramsey numbers: exact values and
  extremal structure}, arXiv:2604.13434v1 (2026). Its edgeless target is
  distinct from our connected claw target; none of its computational
  claims is assumed here.
\end{enumerate}

\end{document}
